\section*{Projects}
Selected work in urban science, architecture, and urban design. For a complete list of projects, please see the \textbf{Projects} section of my website and portfolio: \href{https://arielnoyman.com/}{arielnoyman.com}

\subsection*{MIT Media Lab}

\begin{tablist}

    \item[2015--present] \tab \textbf{City Science Network:} {Helped establish a global network of City Science Labs, including Andorra (2015), Hamburg (2015), Helsinki (2017), Shanghai (2017), Detroit (2019), Ho Chi Minh City (2020), and the Negev (2022).}
    \begin{tablist}
        \item[2022--present] \tab {\textbf{Negev Urban Research (NUR)}, an MIT/Israeli Government urban science lab in the Negev region. Initiated the lab, assembled the team, and secured \$2.5M in government funding.}
        \item[2015--present] \tab {\textbf{HCU CityScience Lab} in Hamburg, Germany -- a multi-year cooperation agreement between MIT, the Free and Hanseatic City of Hamburg, and HafenCity Universität Hamburg. Coordinated MIT--HCU research and projects; helped secure over 5M EUR in funding.}
    \end{tablist}

    \item[2014--present] \tab \textbf{MIT CityScope:} {An urban modeling, simulation, and decision-making platform developed at the MIT City Science Center and deployed internationally.}
    \begin{tablist}
        \item[2016] \tab \href{https://obamawhitehouse.archives.gov/blog/2016/02/23/pcast-releases-technology-and-future-cities-report-president}{President Obama Council of Advisors on Science and Technology (PCAST), Technology and the Future of Cities Report, p. 24}
    \end{tablist}

    \item[2026] \tab \textbf{Beyond the Rulebook:} {Agent-based and cognitive-network spatial simulations in GAMA exploring model-guided pedestrian decisions under uncertainty across subway station and MIT campus accessibility studies.}
    \begin{tablist}
        \item[2026] \tab \href{https://www.arielnoyman.com/projects/26agents/}{Project Documentation \& Recorded Simulations}
    \end{tablist}

    \item[2026] \tab \textbf{Gridnberg:} {Topography-aware pedestrian routing dataset and interactive web application for New York City enriching over 313,000 sidewalk segments with elevation and slope-sensitive accessibility costs.}
    \begin{tablist}
        \item[2026] \tab \href{https://www.arielnoyman.com/gridnberg/}{Gridnberg Interactive Routing Web App}
    \end{tablist}

    \item[2024--2026] \tab \textbf{The Right to Shelter:} {Spatial analysis and collaborative planning tools for expanding bomb-shelter access in marginalized communities using spatial data analysis and machine learning methods.}
    \begin{tablist}
        \item[2025] \tab \href{https://www.ynetnews.com/real-estate/article/s1zbwjtyxx}{\enquote{Innovative algorithm aims to protect Negev Bedouin villages from missile attacks}}
    \end{tablist}

    \item[2024--2025] \tab \textbf{TravelAgent:} {Generative-agent simulation platform for modeling pedestrian wayfinding, spatial experience, and human-like decision-making in 3D urban environments.}
    \begin{tablist}
        \item[2025] \tab \href{https://www.arielnoyman.com/travelAgent/}{TravelAgent Project Page \& Interactive Demos}
    \end{tablist}

    \item[2018--2021] \tab \textbf{CityScope Grasbrook:} {Multi-participant platform for design collaboration and evaluation of proposals for HCH and Herzog \& de Meuron (HdM).}

    \item[2015--2016] \tab \textbf{CityScope FindingPlaces:} {HCI platform for public participation in refugees' accommodation planning process.}
    \begin{tablist}
        \item[2019] \tab \href{https://www.youtube.com/watch?v=PZg6A_A65lQ}{FindingPlaces \& CityScienceLab at the \enquote{City Science Summit}, Hamburg, Germany}
        \item[2018] \tab \href{https://datasmart.ash.harvard.edu/news/article/ar-transforming-tech-what-can-it-do-cities}{AR Is Transforming Tech. What Can It Do for Cities?}
        \item[2017] \tab \href{https://www.media.mit.edu/publications/finding-places/}{Finding Places: HCI Platform for Public Participation in Refugees' Accommodation Process}
        \item[2016] \tab \href{https://medium.com/mit-media-lab/shifting-priorities-finding-places-9ad3bdbe38b8}{Shifting Priorities, Finding Places | MIT Media Lab}
    \end{tablist}

    \item[2014--2015] \tab \textbf{CityScope BRT for Boston:} {Interactive urban planning platform for participatory transit design.}
    \begin{tablist}
        \item[2016] \tab \href{https://www.media.mit.edu/projects/CityscopeBostonBRT/overview/}{Boston Bus Rapid Transit (BRT) Community Engagement}
        \item[2016] \tab \href{http://www.bostonbrt.org/mit-nuestra-comunidad-brt-lego-workshops}{MIT + Nuestra Comunidad BRT LEGO Workshops, Boston BRT}
        \item[2015] \tab \href{https://www.bloomberg.com/news/articles/2015-10-16/mit-s-new-interactive-lego-urban-tool-brings-transparency-to-transit-planning}{MIT's New Interactive LEGO Urban Tool Brings Transparency to Transit Planning}
        \item[2015] \tab \href{http://www.brt.cl/new-tools-for-a-new-era-of-open-transportation-planning/}{BRT+ Centre of Excellence}
    \end{tablist}

\end{tablist}


\subsection*{MIT Center for Advanced Urbanism}

\begin{tablist}

    \item[2013--2014] \tab \textbf{Rebuild by Design (RBD):}
    {\enquote{The New Meadowlands} -- Winning entry for \$150M HUD multistage planning and design competition following Hurricane Sandy in the Tri-State area.}

\end{tablist}


\subsection*{Design Projects}

\begin{tablist}

    \item[2010--2019] \tab \textbf{National Institute for Biotechnology in the Negev (NIBN)} with Chyutin Architects, Ben-Gurion University, Be'er Sheva, Israel. \$22M academic biotechnology laboratory building. Winner of the 2018 Israeli Design Award (Ot HaItzuv).

    \item[2010] \tab \textbf{The Museum of Tolerance, Jerusalem, Israel} with Chyutin Architects. Winning entry for an international competition and design towards construction of a \$260M campus, including two museum exhibitions, a children's museum, conference center, a 500-seat banquet hall, an 800-seat lecture hall, a 1,000-seat outdoor amphitheater, and two viewing theaters.

    \item[2009] \tab \textbf{The Israel Museum, Jerusalem} with EK-A.
    \begin{tablist}
        \item[\textit{EK-A}] \href{https://www.efrat-kowalsky.co.il/projects-repeater/the-israel-museum%2C-jerusalem}{\enquote{The Israel Museum, Jerusalem}}
        \item[\textit{ArchDaily}] \href{https://www.archdaily.com/373684/israel-museum-james-carpenter-design-associates}{\enquote{Israel Museum / James Carpenter Design Associates}}
    \end{tablist}

    \item[2008] \tab \textbf{Tel-Aviv City Museum (Beit Ha'ir)} with EK-A. A 5,000 sq ft city museum and cultural center in Tel-Aviv's historic City Hall building.

    \item[2008] \tab \textbf{Ordos 100, Inner Mongolia, China} with EK-A.
    \begin{tablist}
        \item[\textit{New York Times}] \href{https://www.nytimes.com/2008/05/01/garden/01mongolia.html}{\enquote{In Inner Mongolia, Pushing Architecture's Outer Limits}}
        \item[\textit{CNN}] \href{https://www.cnn.com/style/article/china-ordos-ghost-town/index.html}{\enquote{Ordos, the abandoned city that's an architectural marvel}}
        \item[\textit{ArchDaily}] \href{https://www.archdaily.com/32535/ordos-100-45-efrat-kowalsky-architects}{\enquote{ORDOS 100 \#45: Efrat-Kowalsky Architects}}
        \item[\textit{Vancouver Biennale}] \href{https://vimeo.com/136530514}{\enquote{Ordos 100}}
    \end{tablist}

\end{tablist}
